\subsection{欧氏仿射空间中的正交反射}

沿用前一号的记号，并令 E 为一个仿射空间，其平移空间为 V。给 V 上的型 B 便给 E 赋予欧氏空间结构（《代数》，第 IX 章，§6，第 6 号）。

设 H 是 E 的一个非迷向超平面。关于 H 的对称（《代数》，第 IX 章，§6，第 6 号）也称为关于 H 的正交反射；我们常将其记为 \(s_{H}\)。有 \(s_{H}^{2}=1\)，且 \(s_{H}\) 是 E 的唯一一个位移（同上，定义 3），它不同于恒等变换并固定 H 的元素。与 \(s_{H}\) 对应的 V 的自同构，是关于 H 的方向的正交反射（这是 V 的一个非迷向超平面）。

\(x\) 中的每个 \(\mathbf{E}\) 都能唯一写成 \(x = h + v\) 的形式，其中 \(h \in \mathbf{H}\)，\(v \in \mathbf{V}\) 且与 \(\mathbf{H}\) 正交；有

\[
s _ {\mathrm{H}} (h + v) = h - v.
\]

命题 5. 设 H 和 H' 是 E 中两个平行的非迷向超平面。存在唯一的与 H 正交的向量 \(v \in V\)，使 \(H' = H + v\)。位移 \(s_{H'} s_{H}\) 是向量 2v 的平移。

v 的存在性和唯一性是立即的。与 \(s_{H'}s_{H}\) 对应的 V 的自同构是恒等变换；因此 \(s_{H'}s_{H}\) 是平移。另一方面，设 \(a \in H'\)；于是 \(a - v \in H\)，并且

\[
s _ {\mathrm {H^ {\prime}}} s _ {\mathrm{H}} (a - v) = s _ {\mathrm {H^ {\prime}}} (a - v) = a + v = (a - v) + 2 v,
\]

这表明 \(s_{H'}s_{H}\) 是向量 2v 的平移。

推论. 设 H 和 H' 是两个不同的平行非迷向超平面。如果 K 的特征为零（分别地，p \textgreater{} 0，且 \(p \neq 2\)），则由 \(s_{H}\) 和 \(s_{H'}\) 生成的 E 的位移群是无限二面体群（分别为阶 2p 的二面体群）。

事实上，根据第 IV 章 §1 第 2 号的命题 2，只需证明 \(s_{\mathrm{H}'}s_{\mathrm{H}}\) 是无限阶（分别地，阶为 \(2p\)），这显然成立。

注. 沿用命题 5 的记号，并另假设 K = R。置 \(s = s_{H}\) 和 \(s' = s_{H'}\)。令 \(H_{n}\) 为超平面 \(H + n.v\)，并令 \(C_{n}\) 为 E 中形如 \(a + \xi.v\) 且满足 \(a \in H\) 和 \(n < \xi < n + 1\) 的点所成的集合。\(C_{n}\) 是连通开集，它们构成 \(E - \bigcup_{n} H_{n}\) 的一个划分。因此，它们是 E 中由系统 \(\mathfrak{H} = (\mathrm{H}_{n})_{n \in \mathbf{Z}}\) 定义的室。平移 \((s's)^{n}\) 将室 \(C = C_{0}\) 变换为室 \(C_{2n}\)，且由于 \(s(\mathrm{C}_{0}) = \mathrm{C}_{-1}\)，有 \((s's)^{n}s(\mathrm{C}) = \mathrm{C}_{2n-1}\)。由此，由 s 和 \(s'\) 生成的二面体群 W 简单传递地置换室 \(C_{n}\)。此外，正如下面将要证明的，若室 C 与 w(C) 位于 H 的相反两侧（其中 \(w \in W\)），则 \(l(sw) = l(w) - 1\)（长度是相对于集合 S = \{s, s'\}（《代数》第 IV 章 §1 第 1 号）取的）。事实上，此时有 \(w(\mathrm{C}) = C_{n}\)，其中 n \textless{} 0。若 \(w = (ss')^{k}\)，则 \(sw = (s's)^{k-1}s'\)，所以 \(l(w) = 2k\) 且 \(l(sw) = 2k - 1\)（《代数》第 IV 章 §1 第 2 号，注）。若 \(w = (ss')^{k}s\)，则 \(sw = (s's)^{k}\)，所以 \(l(w) = 2k + 1\) 且 \(l(sw) = 2k\)。

